\IEEEPARstart{B}{efore} a passenger ship sails, its technical staff decide how
many spares of each power-electronic unit to embark: drivers for cabin and
public-area lighting, supplies for charging outlets, fan-coil drives and drives
in the technical spaces. Until the next port where ordered spares can be
received, a failed unit can be replaced only from what is on board, and a spare
of one reference does not replace a failure of another when rating, interface
or mounting differ. The decision is a vector of stocks, judged by whether
\emph{every} reference lasts until resupply; it depends on installed
populations, powered hours, a horizon set by the next real replenishment
rather than by the cruise, the ship's failure history, and supply events that
fail several references at once.

Multi-item spares models cover a system rather than an
item~\cite{sherbrooke2004,basten2014}, including ships that cannot be
resupplied during a mission~\cite{eruguz2018} and kits whose items must all be
available~\cite{mamer1982}, with the demand law given and items usually
independent. Estimated rates need a predictive
distribution~\cite{prak2017}, gamma mixing is the usual treatment of an
uncertain rate~\cite{aronis2004}, and random effects are standard for field
data~\cite{meeker2022}. Left open is how the \emph{structure} of fleet
heterogeneity enters the joint decision: spread shared by all references of a
ship and spread specific to one reference on one ship can give the same
marginal laws but different dependence between references.

The central result is a provisioning guarantee for that unknown split. With
the marginal laws and the total latent variance fixed, the supermodular
ordering of Gaussian vectors~\cite{muller2000} makes the joint cover of every
stock vector non-decreasing in the ship-wide share
(Proposition~\ref{prop:mono}). The ordering is known; this paper applies it
to joint provisioning, draws the consequence that the optimum for independent
references is the cheapest vector feasible for every split, and evaluates what
that guarantee costs and omits on board. Four kinds of evidence are kept
apart: \emph{mathematical} (the proposition and certified optima under known
parameters, among them a large passenger ship); \emph{simulated} (a bridge to
predictive draws, the identification of the two components, and a fleet built
from public data that separates learning field rates, estimating each ship's
marginals and modelling dependence); \emph{public naval data} (ship sizes,
datasheet rates and casualty counts, for orders of magnitude only); and
\emph{terrestrial observation} (street-lighting reports, which illustrate the
grouping).

No public record of shipboard failures per unit with exposure was found for
these products. The Japanese ship-equipment reliability data collected from
1982 were published only as aggregate rates for equipment classes of merchant
ships, about 2.1 failures of electrical and electronic equipment per 1\,000
hours of equipment existence time, corrected for gaps in each ship's
reporting~\cite{kiriya2001}; casualty
databases count losses of electrical power of whole ships without exposure or
damage to units~\cite{emsa2025,maib_portal}; reliability handbooks give
prediction factors, not field rates~\cite{milhdbk217f}. None identifies
ship-to-ship or ship--reference variation, so every shipboard result below on
failures, cover or stock is mathematical or simulated, with published sizes
and specifications as inputs.
